\documentclass[10pt,a4paper]{amsart}
\usepackage[margin=19mm]{geometry}
\usepackage{amsmath,amssymb,mathtools}
\usepackage[scr=rsfs,cal=euler]{mathalfa}
\usepackage{xfrac}
\usepackage[unicode,hidelinks,bookmarks=false]{hyperref}
\newcommand{\ie}{i.e.\ }
\newcommand{\cf}{cf.\ }
\newcommand{\resp}{respectively\ }
\newcommand{\Subsection}{\subsection}
\providecommand{\nobreakdash}{\nobreak\hbox{-}}
\setlength{\emergencystretch}{2em}
\multlinegap0pt
\usepackage{amssymb}
\usepackage{enumerate}
\usepackage[shortlabels]{enumitem}
\usepackage{float}
\usepackage{xcolor}
\newtheorem{assumption}{Assumption}
\newtheorem{lemma}{Lemma}
\newtheorem{theorem}{Theorem}
\newcommand{\C}{\mathbb{C}}
\newcommand{\E}{\mathbb{E}}
\newcommand{\Q}{\mathbb{Q}}
\newcommand{\R}{\mathbb{R}}
\newcommand{\Rr}{\mathcal{R}}
\newcommand{\dd}{\mathrm{d}}
\title{Variance-Optimal Hedging in the Rough Hawkes--Heston Model}
\author{Yingli Wang, Xiaoyu Wang}
\date{Source: arXiv:2609.08541v1 (CC BY 4.0)}
\begin{document}
\maketitle

\noindent\emph{Source and licence.} Variance-Optimal Hedging in the Rough Hawkes--Heston Model. Version arXiv:2609.08541v1 (CC BY 4.0), distributed by arXiv under the Creative Commons Attribution 4.0 International licence, http://creativecommons.org/licenses/by/4.0/, as recorded in the arXiv licence metadata for this exact version and shown on the abstract page https://arxiv.org/abs/2609.08541v1. Full licence notice: \texttt{licenses/arxiv-2609.08541-CC-BY-4.0.txt}.\par\smallskip

\noindent\textit{Editorial scope.} Counted unit: the article's theorem thm:original-compact-hedge (source0.tex lines 1443--1458). The article's complete original proof of it is delivered with it (source0.tex lines 1460--1535, one \texttt{proof} environment of 76 lines). No mathematics is written, repaired or re-derived: the statement and the proof are the article's own lines, in the article's own notation, and no retained line is substituted or re-flowed.  Retained uncounted as the source-local context the proof needs: lines 372--387; lines 415--428; lines 573--577; lines 1252--1259; lines 1275--1303; lines 1407--1410; lines 1430--1438.  Nothing was omitted from any retained span, so the package carries no omission record.  No retained line contains a citation, so the wrapper carries no bibliography and no bibliography entry.  No retained line contains a figure environment, an image-inclusion command or drawing code, so no image is delivered and the package declares no asset dependency.  Editorial formatting only, in addition: an AMS \texttt{amsart} preamble that restates the article's own declarations and macros used by the retained lines, namely the environments \texttt{assumption}, \texttt{lemma}, \texttt{theorem} on separate counters, together with the standard packages those lines need.  The retained environments print in this document as Assumption 1, Lemma 1, Theorem 1 in order of appearance, whereas the article prints the same items with the numbering of its own full document.  The complete native source of the article as distributed in the arXiv e-print for this exact version is bundled unchanged as \texttt{sources/209723/source0.tex}.
\begin{assumption}[Model admissibility]
\label{ass:model}
The kernel $K\in L^2_{\mathrm{loc}}(\R_+)$ is nonnegative, nonincreasing,
not identically zero, and continuously differentiable on $(0,\infty)$. It has
a resolvent of the first kind $L$ which is nonnegative and nonincreasing in
the sense that $s\mapsto L([s,s+h])$ is nonincreasing for every $h\geq0$.
The deterministic curve $g_0$ is continuous and nondecreasing with
$g_0(0)\geq0$, and $\nu$ is a nonnegative Borel measure on $(0,\infty)$ such
that
\[
 \int_{(0,\infty)}z^2\nu(\dd z)<\infty.
\]
Finally, a nonnegative predictable process $V$ with paths in
$L^2_{\mathrm{loc}}(\R_+)$ is given which solves the stochastic Volterra
equation below in the weak sense.
\end{assumption}

\begin{assumption}[Conditional-transform regularity]
\label{ass:transform}
The first-kind resolvent has the decomposition
\[
 L(\dd t)=\ell(t)\dd t+\ell_0\delta_0(\dd t),
 \qquad \ell\in L^1_{\mathrm{loc}}(\R_+),\quad \ell_0\geq0.
\]
Moreover, writing $\Delta_\varepsilon K(t)=K(t+\varepsilon)$,
\[
 \sup_{\varepsilon\in(0,T]}
 \|\Delta_\varepsilon K*L\|_{\mathrm{TV}([0,T])}<\infty
 \qquad\text{for every }T>0.
\]
\end{assumption}

\begin{assumption}[Square-integrability]
\label{ass:l2}
The stock $S$ is a square-integrable $\Q$-martingale on $[0,T]$,
and every target payoff $Y$ under consideration belongs to $L^2(\Q)$.
\end{assumption}

\begin{assumption}[Kernel regularization]
\label{ass:original-kernel-regularization}
For every $n$, the kernel $K_n$ satisfies the kernel conditions in
Assumption~\ref{ass:model}, and
\[
 \epsilon_n:=\|K_n-K\|_{L^1(0,T)}\longrightarrow0.
\]
\end{assumption}

\begin{lemma}[Riccati and memory approximation]
\label{lem:original-deterministic-stability}
For every $R<\infty$, there is a constant $C_R<\infty$ such that, for all
sufficiently large $n$,
\begin{equation}
\label{eq:original-riccati-stability}
 \sup_{|\lambda|\leq R}\sup_{0\leq s\leq T}
 |\psi_{w_\lambda}^n(s)-\psi_{w_\lambda}(s)|\leq C_R\epsilon_n.
\end{equation}
The solutions are uniformly bounded on this compact Fourier interval.
Set $F_n^w(s)=\Rr(w,\psi_w^n(s))$ and $F^w(s)=\Rr(w,\psi_w(s))$, and define
\begin{align}
 a_n^w(t,r)&=\int_t^T F_n^w(T-s)K_n(s-r)\dd s,
       &&0\leq r\leq t\leq T,\label{eq:original-memory-coefficient}\\
 U_n^w(t)&=\int_t^T F_n^w(T-s)g_0(s)\dd s
       +\int_{[0,t)}a_n^w(t,r)\dd Z_r.
       \label{eq:original-log-memory}
\end{align}
Let $a^w,U^w$ denote the same expressions with $K,\psi_w$. Then
\begin{equation}
\label{eq:original-memory-stability}
 \sup_{|\lambda|\leq R}\sup_{r\leq t}|a_n^{w_\lambda}(t,r)-a^{w_\lambda}(t,r)|
 \leq C_R\epsilon_n,
 \qquad
 \sup_{|\lambda|\leq R}\E^\Q\!\int_0^T
 |U_n^{w_\lambda}(t)-U^{w_\lambda}(t)|\dd t
 \leq C_R\epsilon_n.
\end{equation}
\end{lemma}

\begin{equation}
\label{eq:original-transform-envelope}
 |H_{t-}(w_\lambda,T)|^2\leq S_{t-}/S_0.
\end{equation}

\begin{align}
 \beta_n^w(t)&=\frac{1}{D}\left[
 w+\rho\sqrt c\,\psi_w^n(T-t)
 +\int_{(0,\infty)}\left(e^{(\psi_w^n(T-t)-\Lambda w)z}-1\right)q(z)\nu(\dd z)
 \right],\label{eq:original-approx-beta}\\
 \widetilde\vartheta_n^w(t)&=
 \frac{\widetilde H_n^w(t)}{S_{t-}}\beta_n^w(t).
 \label{eq:original-approx-holding}
\end{align}

\begin{theorem}[Hedge convergence on compact Fourier intervals]
\label{thm:original-compact-hedge}
Suppose Assumptions~\ref{ass:model}, \ref{ass:transform}, \ref{ass:l2},
and~\ref{ass:original-kernel-regularization} hold. For each $R<\infty$,
the predictable strategies~\eqref{eq:original-approx-holding} belong to
$L^2(S)$ and
\begin{equation}
\label{eq:original-compact-convergence}
 \sup_{|\lambda|\leq R}
 \|\widetilde\vartheta_n^{w_\lambda}
       -\vartheta^{w_\lambda}\|_{L^2(S)}\longrightarrow0.
\end{equation}
Consequently, for every $a\in L^1([-R,R];\C)$, the corresponding Fourier
integrals of the approximate holdings converge in $L^2(S)$ to the
Fourier integral of the exact holdings.
\end{theorem}

\begin{proof}
By the stock bracket formula and square-integrability,
$\E^\Q\int_0^T S_{t-}^2V_t\dd t<\infty$. Thus
\begin{equation}
\label{eq:original-weight-integrability}
 \E^\Q\int_0^T S_{t-}V_t\dd t
 \leq\left(\E^\Q\int_0^T S_{t-}^2V_t\dd t\right)^{1/2}
      \left(\E^\Q\int_0^T V_t\dd t\right)^{1/2}<\infty.
\end{equation}
For each fixed $n$ and $w$, the deterministic coefficient $\beta_n^w$ is
bounded on $[0,T]$. Setting $W_t=S_{t-}V_t/S_0$, the envelope gives
\[
 \|\widetilde\vartheta_n^w\|_{L^2(S)}^2
 =D\E^\Q\int_0^T V_t|\beta_n^w(t)\widetilde H_n^w(t)|^2\dd t
 \leq D\|\beta_n^w\|_\infty^2\E^\Q\int_0^T W_t\dd t<\infty.
\]
Thus admissibility follows from weighted integrability of the random
envelope.

The map $\chi$ is globally $1$-Lipschitz and has modulus at most one:
on each open half-plane its real differential has operator norm at most
one, and $\chi$ is continuous across the imaginary axis. Integrating along
line segments therefore gives the Lipschitz bound.
Together with~\eqref{eq:original-transform-envelope}, this gives
\[
 |\widetilde H_n^w(t)-H_{t-}(w,T)|^2
 \leq\frac{S_{t-}}{S_0}
        \min\{4,|U_n^w(t)-U^w(t)|^2\}.
\]
For uniform convergence in $|\lambda|\leq R$, define the weight tail
$\mathcal T_W(B):=\E^\Q\int_0^T W_t\mathbf1_{\{W_t>B\}}\dd t$.
For every $B>0$,
\begin{align*}
 \E^\Q\int_0^T V_t|\widetilde H_n^w(t)-H_{t-}(w,T)|^2\dd t
 &\leq 4\mathcal T_W(B)
       +2B\E^\Q\int_0^T|U_n^w(t)-U^w(t)|\dd t.
\end{align*}
Lemma~\ref{lem:original-deterministic-stability} bounds the second term
uniformly by $C_R B\epsilon_n$, while
$\mathcal T_W(B)\to0$ as $B\to\infty$ by~\eqref{eq:original-weight-integrability}.

The deterministic coefficients $\beta_n^w$ are uniformly bounded on
compact Fourier intervals and
\[
 |\beta_n^w(t)-\beta_t^w|
 \leq\frac{|\rho|\sqrt c+\|q\|_{L^2(\nu)}\sqrt{m_2}}{D}
       |\psi_w^n(T-t)-\psi_w(T-t)|.
\]
Indeed, the difference of the two jump exponentials is bounded by
$z|\psi_w^n-\psi_w|$. Finally,
\begin{align*}
 \|\widetilde\vartheta_n^w-\vartheta^w\|_{L^2(S)}^2
 &=D\E^\Q\int_0^T V_t
       |\beta_n^w(t)\widetilde H_n^w(t)-\beta_t^wH_{t-}(w,T)|^2\dd t\\
 &\leq 2D\|\beta_n^w\|_\infty^2
       \E^\Q\int_0^T V_t|\widetilde H_n^w(t)-H_{t-}(w,T)|^2\dd t\\
 &\quad+2D\|\beta_n^w-\beta^w\|_\infty^2
       \E^\Q\int_0^T W_t\dd t.
\end{align*}
Consequently, for all sufficiently large $n$ and every $B>0$,
\begin{equation}
\label{eq:original-weighted-error-bound}
 \sup_{|\lambda|\leq R}
 \|\widetilde\vartheta_n^{w_\lambda}-\vartheta^{w_\lambda}\|_{L^2(S)}^2
 \leq C_R\bigl(\mathcal T_W(B)+B\epsilon_n+\epsilon_n^2\bigr),
\end{equation}
where $C_R$ is independent of $n$ and $B$. Taking $n\to\infty$ at fixed
$B$, then $B\to\infty$, proves~\eqref{eq:original-compact-convergence}.
For fixed $n$, continuous dependence of $\psi_w^n$ and
$F_n^w$ on $\lambda$, followed by the same memory and weighted truncation
estimates, gives continuity of
$\lambda\mapsto\widetilde\vartheta_n^{w_\lambda}$ in $L^2(S)$; the exact
holdings have the same property. This ensures strong measurability of the
Fourier integrands. Minkowski's inequality then
proves the final assertion.
\end{proof}

\end{document}
